% Paper C: object-grid wafer classification / categorical-resize rule
% Skeleton (article style; retarget to IEEE T-Semi (IEEEtran) at submission).
% Prose source: ../DRAFT.md
\documentclass[11pt]{article}
\usepackage[margin=1in]{geometry}
\usepackage{booktabs}
\usepackage{amsmath,amssymb}
\usepackage{graphicx}
\usepackage{hyperref}
\usepackage{xcolor}
\newcommand{\todo}[1]{\textcolor{red}{[TODO: #1]}}

\title{A 1.16M-Parameter Object-Grid Network Beats an 88M Pixel CNN for
Wafer Map Classification:\\ Categorical Signals Must Not Be Interpolated}
\author{\todo{authors}}
\date{}

\begin{document}
\maketitle

\begin{abstract}
Wafer-map defect classification is dominated by large pixel CNNs applied
to high-resolution fail-bit images. We show that a two-stage object-aware
pipeline --- a chip-level classifier that labels each die region with an
object type, followed by a 1.16M-parameter CNN over the native
$32{\times}32$ chip-grid of one-hot object labels --- reaches val macro-F1
$0.9946$ on a 33-class wafer benchmark, beating an 88M ImageNet-pretrained
ConvNeXtV2 on raw pixels ($0.9851$) with a $76\times$ smaller model and
${\sim}750\times$ faster training, and exceeding the $0.9919$ oracle
ensemble ceiling of the constituent models. The enabling factor is
representational: object-identity maps are \emph{categorical}, and any
interpolating resize corrupts them --- the same information presented as
an interpolated 384\,px compound image caps at $0.9784$, and bicubic
resizing manufactures fractional object identities on $6.5\%$ of pixels of
a typical map. Block-integer expansion or native-resolution processing
removes this ceiling. The pipeline is robust to $10\%$ chip-classifier
error ($0.9917$ at full data) and, counter-intuitively, drops the raw
pixel channel: re-presenting pixels after stage 1 hurts ($0.9707$ vs
$0.9946$). A 10-seed experiment on real public WM-811K maps supports the
categorical principle ($+0.060$ macro-F1 over interpolated-grayscale
input).
\end{abstract}

\section{Introduction}\label{sec:intro}
\todo{port DRAFT.md Sec 1 prose}

\section{Method}\label{sec:method}
\todo{port DRAFT.md Sec 2: stage-1 inline-crop chip CNN; 6-conv ChipGridCNN;
block\_expand rule}

\section{Experiments}\label{sec:experiments}

\begin{table}[t]
\centering
\caption{Fair-evaluation comparison (same split/protocol; measured).}
\label{tab:main}
\begin{tabular}{lllcc}
\toprule
Model & Input & Params & val\_f1 & test\_f1 \\
\midrule
R-only ConvNeXtV2 (pretrained) & 1024\,px pixels & 88M & 0.9851 & --- \\
Compound 3ch BICUBIC 384       & interpolated    & 88M & 0.9784 & 0.9736 \\
Obj-only small CNN             & $32^2$ one-hot  & 0.4M & 0.9844 & --- \\
V3 + pixels (6ch)              & grid + pixels   & 1.16M & 0.9707 & --- \\
\textbf{V3 obj-only (ours)}    & $32^2{\times}5$ one-hot & \textbf{1.16M} & \textbf{0.9946} & \textbf{0.9872} \\
Oracle ceiling (either-correct)& ---             & ---  & 0.9919 & --- \\
\bottomrule
\end{tabular}
\end{table}

\begin{table}[t]
\centering
\caption{Encoding ablation ($n{=}100$/class, seed 42; measured).}
\label{tab:encoding}
\begin{tabular}{llccc}
\toprule
Variant & Encoding & in\_ch & val\_f1 & test\_f1 \\
\midrule
V0 & fail-bit pixels only  & 1 & 43.59\% & 43.85\% \\
V2 & single-object binary  & 2 & 65.43\% & 64.79\% \\
V1 & integer id $/5$       & 2 & 95.05\% & 97.26\% \\
V3 & one-hot 5ch           & 6 & \textbf{96.89\%} & \textbf{98.79\%} \\
\bottomrule
\end{tabular}
\end{table}

\todo{port DRAFT.md Sec 3 incl. 3.1 WM-811K real anchor (10 seeds:
nearest+one-hot 0.7633 $>$ nearest-gray 0.7248 $>$ bicubic-gray 0.7036)}

\section{Discussion}\label{sec:discussion}
\todo{port DRAFT.md Sec 4 prose (sufficiency, pixels-as-noise, operations,
composition with companion work)}

\section{Conclusion}\label{sec:conclusion}
\todo{write from DRAFT.md headline claims}

\bibliographystyle{IEEEtran}
% \bibliography{refs}
\todo{refs.bib per DRAFT.md Sec 5 (verify WM-811K Wu2015, MixedWM38
Wang2020, ConvNeXtV2 Woo2023, R-CNN)}

\end{document}
